% =============================================================================
%  capitulos/cap1_introducao.tex
%  CAPÍTULO 1 — INTRODUÇÃO
%
%  Objetivo: contextualizar o leitor sobre o trabalho.
%  Responde: O que é? Por quê foi desenvolvido? Para quê serve?
%            Como o trabalho está organizado?
% =============================================================================

\chapter{Introdução}
\label{cap:introducao}
% \label permite referenciar este capítulo em qualquer outro arquivo com
% \cref{cap:introducao}, que gera "Capítulo 1" automaticamente.


% -----------------------------------------------------------------------------
%  1.1 CONTEXTUALIZAÇÃO
%  Apresente o cenário atual da área onde o problema se insere.
%  Use dados e referências para mostrar que o problema é real e relevante.
%  Responde: "Em que mundo este software existe?"
% -----------------------------------------------------------------------------
\section{Contextualização}
\label{sec:contextualizacao}

TODO: Descreva o cenário atual da área. Apoie-se em dados, estatísticas e
referências bibliográficas. Exemplo: segundo \citeonline{pressman2016}, a
demanda por sistemas X cresceu Y\% nos últimos anos. O setor de Z ainda
depende de processos manuais, o que gera impactos como...


% -----------------------------------------------------------------------------
%  1.2 PROBLEMA
%  Delimite a lacuna ou dificuldade que o software resolve.
%  Deve ser objetivo: no máximo dois parágrafos.
%  Responde: "Qual dor exatamente você está atacando?"
% -----------------------------------------------------------------------------
\section{Problema}
\label{sec:problema}

TODO: Descreva de forma clara e direta o problema central. Exemplo:
``A ausência de uma ferramenta integrada para X resulta em Y, gerando
impactos negativos como Z para os usuários de W.''


% -----------------------------------------------------------------------------
%  1.3 OBJETIVOS
% -----------------------------------------------------------------------------
\section{Objetivos}
\label{sec:objetivos}

% --- Objetivo Geral ---
% Uma frase. O que você entrega ao final do trabalho?
% Use verbo no infinitivo: Desenvolver / Criar / Implementar / Propor...
\subsection{Objetivo Geral}
\label{subsec:objetivo-geral}

TODO: Desenvolver [o quê] para [quem / qual finalidade], utilizando
[abordagem tecnológica principal].


% --- Objetivos Específicos ---
% Os passos concretos para atingir o objetivo geral.
% Cada item deve ser verificável — você vai retomar esses objetivos na Conclusão.
% Use verbos no infinitivo: Analisar, Modelar, Implementar, Validar, Documentar...
\subsection{Objetivos Específicos}
\label{subsec:objetivos-especificos}

Para alcançar o objetivo geral, foram definidos os seguintes objetivos
específicos:

\begin{itemize}
  \item TODO: Levantar e analisar os requisitos funcionais e não funcionais
        do sistema junto aos usuários e partes interessadas;
  \item TODO: Projetar a arquitetura do sistema e modelar os dados;
  \item TODO: Implementar os módulos de [funcionalidade principal];
  \item TODO: Realizar testes de [tipo] para validar o funcionamento do
        sistema frente aos requisitos levantados;
  \item TODO: Adicionar demais objetivos específicos do seu projeto.
\end{itemize}


% -----------------------------------------------------------------------------
%  1.4 JUSTIFICATIVA
%  Por que é importante desenvolver esta ferramenta?
%  Aborde impacto: social, técnico, acadêmico ou econômico.
%  Responde: "Por que o mundo fica melhor com este software existindo?"
% -----------------------------------------------------------------------------
\section{Justificativa}
\label{sec:justificativa}

TODO: Argumente a relevância do trabalho. Considere: quem se beneficia
diretamente, qual o impacto esperado e por que uma solução de software é a
abordagem mais adequada para este problema.


% -----------------------------------------------------------------------------
%  1.5 ORGANIZAÇÃO DO TRABALHO
%  Guia rápido para o leitor saber o que encontrará nos próximos capítulos.
%  Mantenha um parágrafo simples — sem exagerar no detalhamento.
% -----------------------------------------------------------------------------
\section{Organização do Trabalho}
\label{sec:organizacao}

Este trabalho está organizado da seguinte forma: o
\cref{cap:fundamentacao} apresenta a fundamentação teórica, abordando as
tecnologias e padrões adotados no desenvolvimento. O
\cref{cap:metodologia} descreve a metodologia ágil utilizada e a
organização das etapas do projeto. O \cref{cap:requisitos} detalha os
requisitos funcionais, não funcionais e a modelagem UML do sistema. O
\cref{cap:implementacao} descreve a implementação, destacando a
arquitetura, a modelagem de dados e trechos relevantes do código. O
\cref{cap:testes} apresenta os testes realizados e os resultados obtidos.
Por fim, o \cref{cap:conclusao} retoma os objetivos, discute as
dificuldades encontradas e propõe trabalhos futuros.
